\documentclass[a4paper,12pt]{article}
\usepackage[utf8]{inputenc}
\usepackage[T1]{fontenc}
\usepackage[ngerman]{babel}
\usepackage{amsmath}
\usepackage{amssymb}
\usepackage{enumitem}
\usepackage{geometry}
\geometry{a4paper, margin=2.5cm}

% Einheitliche Fall-Überschrift mit Abstand davor und danach
\newcommand{\fall}[1]{\par\medskip\noindent\underline{#1}\par\nopagebreak\smallskip}

\title{Einsendeaufgaben -- Lektion 5\\[0.5em]
\large Modul 61111: Mathematische Grundlagen}
\author{}
\date{}

\begin{document}
\maketitle

\section*{Aufgabe 5.1}

Sei $x \in \mathbb{R}$ beliebig.

\fall{Fall 1: $x < -1$}

Sei $(x_n)$ eine beliebige Folge mit Grenzwert $x$.
Wir definieren ein $\varepsilon := \dfrac{|x - (-1)|}{2}$.

Es gilt dann $x < x + \varepsilon < -1$ und für fast alle $x_n$ gilt $x_n < x + \varepsilon < -1$.
Und daher für fast alle $x_n$ gilt:
\[
    f(x_n) = x_n + 2.
\]

Da $x + 2$ stetig auf ganz $\mathbb{R}$ ist, konvergiert $(f(x_n))$ gegen $f(x)$.

\fall{Fall 2: $x > -1$}

Gleiche Argumentation wie Fall 1.

\fall{Fall 3: $x = -1$}

Wir verwenden das $\varepsilon$-$\delta$-Kriterium.
Sei $\varepsilon > 0$ beliebig.

Da jede Polynomfunktion stetig ist, existiert ein $\delta_1 > 0$, sodass
für alle $y$ mit $|y + 1| < \delta_1$ gilt:
\[
    |(y + 2) - (x + 2)| < \varepsilon.
\]

Genauso existiert ein $\delta_2 > 0$, sodass
für alle $y$ mit $|y + 1| < \delta_2$ gilt:
\[
    |y^2 - x^2| < \varepsilon.
\]

Wir definieren ein $\delta := \min(\delta_1, \delta_2)$.

Zu zeigen ist nun: für alle $y$ mit $|y + 1| < \delta$ gilt:
\[
    |f(y) - f(x)| < \varepsilon.
\]

\fall{Fall 3.1: $y \leq -1$}

Da $\delta \leq \delta_1$ gilt, folgt $|y + 1| < \delta \leq \delta_1$.
Wegen $y \leq -1$ gilt $f(y) = y + 2$.
Aus $|y + 1| < \delta_1$ und $f(y) = y + 2$ folgt:
\[
    |f(y) - f(x)| = |(y + 2) - (x + 2)| < \varepsilon.
\]

\fall{Fall 3.2: $y > -1$}

Gleiche Argumentation.

\end{document}
